% Human civilisation has used microbes and cell cultures since at least the Neolithic for fermentation \cite{liu2018fermented}. 
After the introduction of recombinant protein production in 1977 \cite{itakura1977expression}, a large number of monoclonal antibodies, hormones, enzymes and other proteins of pharmaceutical, industrial and scientific importance are being synthesised using microbes such as bacteria and yeast, insect cells and mammalian cells. Consequently, recombinant proteins currently has a market value in billions of dollars, making it one of the highly valued technologies \cite{walsh2014biopharmaceutical, puetz2019recombinant}. 


There has been much research to improve the recombinant protein production technology. In particular, the development of the pET vector system in 1991 has revolutionised the use of \textit{Escherichia coli} for protein production \cite{dubendorf1991controlling}. \textit{E. coli} is often the host of choice because it is relatively inexpensive, and has a faster growth rate than other expression hosts \cite{Rosano2014-oq, demain2009production}. Now there are a multitude of optimised vectors such as pGEX, pMAL and pET as well as a number of engineered strains of \textit{E. coli} such as BL21 and BL21(D3). Several guidelines and practices have also been proposed to maximise the chances of successful experiments \cite{Berlec2013-mb, Rosano2014-oq}. Furthermore, several high-throughput methods have made the process scalable \cite{stevens2000design, braun2003high, jia2016high}. These new advancements has made the process of recombinant protein production much easier and economical. 


\section{Recombinant protein production}
The initial step in recombinant protein production is a successful protein expression. The expressed protein then needs to be soluble for use in many structural, functional and pharmaceutical studies where concentrated protein samples are desired \cite{Kramer2012-wk, Hou2018-yd}. Despite almost 40 years of refinements in protocols and technology, around half of the recombinant protein expression experiments fail at the expression stage and nearly half of expressed protein are insoluble \cite{targetdbmetrices} (Figure \ref{fig:fail_succ}). This makes the protein production process more challenging. Predicting protein expression and solubility can help plan the experiment and save time and resources. Furthermore, using a highly optimised target gene can increase the success rate of recombinant protein production.


\begin{figure}[htbp!]
\center
\includegraphics[width=1\textwidth]{chapters/Introduction/Figures/Status_of_protein_expression(2).pdf}
\caption[The success rate of recombinant protein production is around a quarter.]{\textbf{The success rate of recombinant protein production is around a quarter. (A)} All experiments, using different vectors and hosts, preformed for deposition to the TargetTrack database shows around $19\%$ of experiments are purified. Data taken from Protein Structural Initiative (PSI:Biology) metrics. \textbf{(B)} A subset of experiments from TargetTrack database using pET vector and \textit{E. coli} as expression host, shows around $50\%$ of experiments produce soluble proteins. }%the List of Figures because of the *}
\label{fig:fail_succ}
\end{figure}


In the following sub-sections, we will discuss these two steps\textemdash protein expression and solubility in details. Unless otherwise stated, these discussions will refer to prokaryotes, in particular, \textit{E. coli} based systems. 

% Protein expression
\input{chapters/Introduction/ProteinExpression/ProteinExpression}

% Solubility
\subsection{Protein solubility}
% The Gibb's phase rule describes the state of a system at equilibrium and is given by 

% \begin{equation}
%     f + p = c + 2
% \end{equation}

% where $c$ is the number of components, $p$ is the number of phases and $f$ is the degree of freedom. For a two phase system (solid and liquid) at constant pressure and temperature, with two components (solute and solvent for example, protein and water), $f = 0$. This means the amount of protein in solution phase no longer depends upon the total amount of protein in the system and remains constant \cite{arakawa19853}. Thermodynamically, the chemical potential of protein in solid phase and liquid phase are now equal. The physical solubility, thus, can be defined as the concentration of protein in a saturated solution which is in equilibrium with the solid phase. Solubility has units of moles per liter ($m/L$) and often, for convenience, written in grams per liter ($g/L$). 


% This physical solubility of a protein depends on the experimental conditions such as temperature, pH, ionic concentration \cite{Kramer2012-wk}. This might lead to the variability of results across various experimental conditions. Hence, 
Solubility is defined as the proportion of the supernatant fraction, obtained after the centrifuging the the translation mixture, to the uncentrifuged total protein \cite{Niwa2009-ye}. It ranges from $0\%$ to $130\%$ with solubility less than $30\%$ categorised as aggregation-prone and greater than $70\%$ are highly soluble. Several \textit{intrinsic} features of the protein itself such as molecular weight, flexibility, hydrophobicity, isoelectric point and structural propensities are also known to influence solubility \cite{Wilkinson1991-zp, Chiti2003-zk, Tartaglia2004-wm, Diaz2010-md}. Sometimes these intrinsic features are modified by doing either mutagenesis or truncation, which might assist in improving solubility. Several solubility enhancing tags are also available for example, thioredoxin (TRX), maltose binding protein (MBP), small ubiquitin-related modifier (SUMO) and glutathione S-transferase (GST). Although, the exact mechanisms of how these tags work is still unclear, it is proposed that they might act like a chaperone and assist in correct folding of the target protein or add charges which decreases the overall aggregation propensity \cite{Costa2014-oe}.

\subsubsection{Intrinsic properties of a protein}
Intrinsic properties are derived from different properties of the residues inside the poly-peptide chain. The commonly used intrinsic properties of proteins are described below.

\paragraph{Hydrophobicity}
Water soluble proteins fold such that the hydrophilic parts are exposed and can form hydrogen bonds with the water molecules whereas the hydrophobic parts are buried in the core. This hydrophobic effect is thought to be responsible for protein solubility \cite{tanford1978hydrophobic}. Several scales have been proposed to measure the hydrophobicity of residues \cite{Kyte1982-qn, abraham1987extension, janin1979surface, rose1985hydrophobicity}. However, none of these scales can fully model the full range of behaviour of residues \cite{charton1982structural}. 


We will use Kyte-Doolittle's scale \cite{Kyte1982-qn} for representative purposes. In this scale, residues are given a hydropathy score such that positive scores represent hydrophobicity and negative score represent hydrophilicity. The magnitude of score represents the strength. For example: isoleucine (I) is given 4.5 and is the most hydrophobic residue, where as  arginine (R) is given -4.5 and is the most hydrophilic residue. Using these scores, a hydropathy plot for a given polypeptide can be drawn (Fig. \ref{fig:hydrophobicity_flexibility_plot}). Hydropathy plot can be used to examine the hydrophobicity of protein region of interest. Using hydrophobic effect, we can then infer whether the residue is buried or located on the surface. Furthermore, the overall hydrophobicity of the protein can be determined by averaging the hydropathy scores across the polypeptide chain to obtain the GRand AVerages of hydropathY (GRAVY) score.


\begin{figure}[htbp!]
\center
\includegraphics[width=1\textwidth]{chapters/Introduction/Figures/flex_and_hydrop.png}
\caption[Profiles of KPC1\_DROME (UniProtKB P05130).]{\textbf{Profiles of KPC1\_DROME (UniProtKB P05130)}. 
\textbf{(A)} Flexibility plot generated using normalised B-factors from Vihinen et al. (1994), with a sliding window of 9 residues. For these normalised B-factors, values greater than one are regarded to be flexible and values less than one are rigid. \textbf{(B)} Hydropathy plot generated using Kyte-Doolittle's hydrophobicity scale, with a sliding window of 19 residues. For illustration, the hydropathy and flexibility of residues at around position $440$ to $470$ (shown by dotted box) are positive and rigid. This indicates the presence of an alpha helix which is supported by the actual 3D structure \textbf{(C)}. The coloured helix is the region $440$ to $470$.  }%the List of Figures because of the *}
\label{fig:hydrophobicity_flexibility_plot}
\end{figure}


% \begin{figure}[htbp!]
% \center
% \includegraphics[width=1\textwidth]{chapters/Introduction/Figures/kpc1_drome.png}
% \caption{\textbf{3D structure of KPC1\_DROME (UniProtKB P05130)}. The highlighted helix runs from position $448$ to $467$. The presence of this helix can be inferred by the hydropathy plot (Fig \ref{fig:hydrophobicity_plot}) and the flexibility profile plot (Fig. \ref{fig:flexibility_plot}). }%the List of Figures because of the *}
% \label{fig:3dplot_kpc1_drome}
% \end{figure}


\paragraph{Isoelectric point}
The net charge of a protein at a given pH depends on acid dissociative constant (pKa) of ionisable groups such as amine and carboxyl group. At a certain pH, the amount of negative and positive charge are equal, resulting a zero net charge. This pH is called the isoelectric point (pI). The net charge is positive at pH below pI and negative at pH above pI \cite{shaw2001effect}. Since there are no net charges, the solubility of a protein is minimum at the isoelectric point.

\paragraph{Instability index}
Guruprasad et. al \cite{guruprasad1990correlation} found that the distribution of certain dipeptides on stable and unstable protein is different. Based on 12 unstable and 32 stable proteins, they assigned a weight called dipeptide instability weight value (DIWV) for all dipeptides. The instability index (II) is then given by equation \ref{eqn:instability_index}.

\begin{equation}
    II = \frac{10}{L}\sum_{i=1}^{L-1}DIWV(x_i y_{i+1})
    \label{eqn:instability_index}
\end{equation}

where $L$ is the length of the sequence and $x_i y_{i+1}$ is a dipeptide.

%%%http://pd.chem.ucl.ac.uk/pdnn/refine3/adps.htm
%%Biomolecular Crystallography: Principles, Practice, and Application : page 641

\paragraph{Flexibility}
Protein molecules are dynamic and are inherently flexible due to the motion of the constituent atoms \cite{vihinen1994accuracy, alvarez2014relationship, teilum2009functional}. The structural flexibility of a protein can be inferred by using B-factors \cite{vihinen1994accuracy, Karplus1985-ea, Smith2003-gb}.


The B-factor (Equation \ref{eqn:bfactor}) or temperature factor of the atoms in a crystalline structure is the measure of mean squared displacement vibration around their mean position $(u = \langle (x-x_0)^2 \rangle)$, where $x$ is the displacement of the atom from its mean position $x_0$. B-factor thus reflects the \textit{orderedness} of the crystal lattice and subsequent uncertainty in X-ray scattering structure determination \cite{Schlessinger2005-ps, Carugo2018-ka, Bramer2018-dh}. It has unit of $\AA^2$.


\begin{equation}
    B = 8\pi^2 u
    \label{eqn:bfactor}
\end{equation}

To understand the effect of the B-factor, we define a quantity $f$ called the atomic scattering factor as:

\begin{equation}
    f = \frac{amplitude\ of\ wave\ scattered\ by\ an\ atom}{amplitude\ of\ wave\ scattered\ by\ one\ electron}
\end{equation}

Atomic scattering factor after taking into account of the motion of atom becomes:

\begin{equation}
    f_B = f\cdot e^{-B(sin\ \theta /\lambda)^2}
\end{equation}

where $\theta$ is the Bragg angle and $\lambda$ is the wavelength of the wave. Thus, we see that B-factor attenuates the amplitude of wave scattered by an atom (Figure \ref{fig:bfactors} ).

% \begin{figure}[htbp!]
% \center
% \includegraphics[width=0.5\textwidth]{chapters/Introduction/Figures/bfactors.pdf}
% \caption{\textbf{Attenuation of the incident waves due to increasing B-factor}. As B-factor increases, atomic scattering factor $(f)$ and consequently, the amplitude of wave scattered by an atom decreases rapidly.}%the List of Figures because of the *}
% \label{fig:bfactors}
% \end{figure}

Experimental B-factors for different residues in a protein sequence can be obtained the from Protein Data Bank (PDB). However, due to variation of structures, the B-factor of a given residue varies, even within the same polypeptide chain. For standardisation, the B-factor of residues within a chain are normalised using a z-score.

\begin{equation}
    B_{norm}^i =  \frac{B^i - \langle B\rangle}{\sigma}
\end{equation}

where, $B_{norm}^i$ is the normalised B-factor of residue $i$, $\langle B\rangle$ is the mean and $\sigma$ is the standard deviation of all B-factors across the chain.

\begin{wrapfigure}{r}{0.5\textwidth}
  \begin{center}
    \includegraphics[width=0.45\textwidth]{chapters/Introduction/Figures/bfactors.pdf}
    \caption[Attenuation of the incident waves due to increasing B-factor.]{\textbf{Attenuation of the incident waves due to increasing B-factor.} As B-factor increases, atomic scattering factor $(f)$ and consequently, the amplitude of wave scattered by an atom decreases rapidly.}%the List of Figures because of the *}
    \label{fig:bfactors}
  \end{center}
\end{wrapfigure}


A number of high resolution structures are sampled from database and $B_{norm}^i$ is calculated for each residue on each protein. The mean of $B_{norm}^i$ across the sampled structures is the final normalised B-factor \cite{Schlessinger2005-ps, Smith2003-gb, Karplus1985-ea, vihinen1994accuracy}. 


The structural flexibility of protein can either be determined by using these normalised B-factors from experiments or by directly computing atomic displacements form molecular dynamics simulations (MDS) 
\cite{dong2018structural, kufareva2011methods}. Root mean square fluctuations and radius of gyration from MDS are useful in examining the flexibility. The profile plot (Fig. \ref{fig:hydrophobicity_flexibility_plot}) can be used to visualise and infer the local flexibility and dynamics of the protein structure. Since structural flexibility is inherently related to the protein dynamics, it is thought to influence several properties such as conformal variations, functions, thermal stability, ligand binding and disordered regions  \cite{Vihinen1987-jo, Teague2003-vq, Ma2005-cr, Yuan2005-gl, Yin2011-su, amaral2017protein}. Although the relationship of flexibility with solubility has been noted previously \cite{Tsumoto2003-qp}, it has been overlooked. 


% \begin{figure}[htbp!]
% \center
% \includegraphics[width=1\textwidth]{chapters/Introduction/Figures/flexibility.pdf}
% \caption{\textbf{Flexibility profile of KPC1\_DROME (UniProtKB P05130)}. Plot was generated using normalised B-factors from Vihinen et al. (1994), with a sliding window of 9 residues. For these normalised B-factors, values greater than one are regarded to be flexible and values less than one are rigid. For illustration purposes, the residues at around position $440$ to $470$ are rigid, suggesting that there might be some structures. This is also supported by the hydropathy plot (Fig. \ref{fig:hydrophobicity_plot}) and the actual 3D structure (Fig. \ref{fig:3dplot_kpc1_drome}).}%the List of Figures because of the *}
% \label{fig:flexibility_plot}
% \end{figure}


Beside these features, amino acids also tend to have different structural propensities which is also thought to influence protein solubility \cite{Idicula-Thomas2005-qw, Huang2012-ft}.


Many solubility prediction tools have been developed around these features using statistical models (e.g., linear and logistic regressions) and machine learning models (e.g., support vector machines and neural networks) \cite{Hirose2013-nq, Habibi2014-jq, Hebditch2017-bg, Sormanni2017-lo, Heckmann2018-wb, Wu2019-nz, Yang2019-kd}. Newer tools such as SOLart also employ 3D structural information for a precise estimation of solvent accessibility, which makes the prediction more accurate \cite{hou2020solart}. Despite a higher prediction accuracy, the usability of these structure based tools might be limited due to the lack of 3D structure information of many proteins of interest.





% Signal peptides
\section{Signal Peptide}
Secretory proteins such as hormones and toxins are some of the commercially important use cases of recombinant protein expression. These secretory proteins are often enriched with a short hydrophobic peptide at the N-terminal (Figure \ref{fig:signal_peptides}). This is called signal peptide (SP). Signal peptides are recognised by a protein complex called the signal recognition particle (SRP). SRP carries the signal peptide to the endoplasmic reticulum (ER) lumen, where post translation modification happens and the newly synthesised protein is secreted out.  Despite having no consensus, SP have a tripartite structure as N-region, H-region and C-region \cite{Von_Heijne1985-qv} (Figure \ref{fig:signal_peptide_structure}). 

\begin{figure}[htbp!]
\center
\includegraphics[width=1\textwidth]{chapters/Introduction/SignalPeptide/Figures/signal_peptide.png}
\caption[Tripartite structure of a signal peptide.]{\textbf{Tripartite structure of a signal peptide.} N, H and C domains in a signal peptide. C domain also contains a cleavage site from which mature peptide is cleaved off after reaching the endoplasmic lumen. Sequence logo (bottom) shows an enrichment of Leucine (L) at the H-region.}
\label{fig:signal_peptide_structure}
\end{figure}


\begin{wrapfigure}{r}{0.5\textwidth}
  \begin{center}
    \includegraphics[width=0.45\textwidth]{chapters/Introduction/SignalPeptide/Figures/sp_nosp_lineplot.pdf}
    \caption[Signal peptides are highly hydrophobic at the N-terminal.]{\textbf{Signal peptides are highly hydrophobic at the N-terminal.} Savitzky–Golay filtering is applied to the hydrophobicity. SP (N=2,609)and NO\_SP(N=14,655) sequences from SignalP 5.0 dataset (Armenteros \textit{et al.} (2019)). SP, Signal Peptide; NO\_SP, Not a SP.}%the List of Figures because of the *}
    \label{fig:signal_peptides}
  \end{center}
\end{wrapfigure}
The N-region usually consists of 2-5 charged residues, whereas the H-region is highly hydrophobic and forms an alpha helix. The C-region consists of small polar uncharged residues which often form a $\beta$-sheet structure. This topology is thought to help binding with signal peptidase and cleave off the signal peptide at the cleavage site. 

% Translation is temporarily paused when signal recognition particle (SRP) binds to the SP. SRP then carries the ribosome/mRNA complex towards the translocon of endoplasmic recticulum (ER). Signal peptide is then detached from the mature peptide and possibly digested by signal peptidase. Translation resumes, pushing the nascent polypetide inside the ER lumen, where post translation modification happens and the newly synthesised protein is secreted out. 

The use of signal peptides in recombinant expression can lead to a high yield. As an added benefit, the proteins are often closer to the native activity because proper folding can happen inside the ER \cite{Futatsumori-Sugai2010-iu, Karyolaimos2019-ip}. 



% Optimisation methods
\input{chapters/Introduction/OptimisationMethods/OptimisationMethods}

% Classifier
\section{Classifiers}
Classifiers are functions which map the input vector to some specific category. In the context of machine learning, we usually use a set of labelled data (known as training data) to fit a classifier. This fitted classifier, also known as model, can then be used to do predictions on unknown data. Many types of classifiers exist such as linear classifiers, support vector machines, decision trees and neural networks (Figure \ref{fig:classifier_comparision}). One special classifier built from an ensemble of decision trees called the Random forest classifier is used in this work. 

\subsection{Random forest}
A random forest is a classifier consisting of a collection of tree structured classifiers $\{h(x, \Theta_k), k=1, ...\}$ where the  $\{\Theta_k\}$ are independent, identically distributed random vectors and each tree casts a unit vote for the most popular class at input {x} \cite{breiman2001random}. For a given data, a number of decision trees $B$ are constructed using bootstrapping. Every time a new split is performed, a subset $m$ of total vectors (features) $f$ is used, such that $m < f$. If $m = f$, this procedure is called bagging. In practise, usually, $m \approx \sqrt{f}$. Surprisingly, the number of trees, $B$, is not critical and setting this to a very high value will not lead to overfitting, which can be shown to be a consequence of the strong law of large numbers \cite{james2013introductiontostatlearning, breiman2001random}. This also makes random forests more robust to noise and outliers than other classifiers. 

Typically, each bagged tree uses around two-third of training points. Out-of-bag (OOB) estimates (errors) can be computed by preforming predictions on the remaining one-third points. OOB errors reflect the generalisability of the model. For each feature, we can also compute the total reduction of Gini index by that feature at each tree split, where Gini index is defined by: 
\begin{equation*}
    G = \sum_{k=1}^{K}(p_{mk}(1-p_{mk}))
\end{equation*}
and is a measure of variance across $K$ classes for the proportion of training observations, $p_{mk}$, in the $m^{th}$ region that are from the $k^{th}$ class. This gives us the feature importance, sometimes called the Gini importance. However, it should be noted that $m$ features used for each split are actually random and is not based on feature importance at all. This seemingly deceptive strategy forces all the subsequent trees to not use the same strongest predictor, thus forcing a decorrelation among the trees. This makes the outputs less variable and more reliable \cite{james2013introductiontostatlearning}.

\begin{figure}[!hbtp]
\centerline{\includegraphics[width=1\textwidth]{chapters/Introduction/Figures/classifier_comparision.pdf}}
\caption[Comparison between some of the commonly used classifiers on synthetic datasets..]{{\bf Comparison between the some of the commonly used classifiers on synthetic datasets.} Three commonly used classifiers\textemdash Support Vector Machine (SVM) with Radial Basis Function (RBF) kernel, Random forest and a simple neural network (multi-layered perceptron) on three synthetic datasets as input. The input data is a binary data coloured as red and blue. The decision boundary is shown as contours and the accuracy of each classifier is shown on bottom right. Adapted from https://scikit-learn.org/stable/auto\_examples/classification/plot\_classifier\_comparison.html.}\label{fig:classifier_comparision}
\end{figure}






